\chapter{Abdominal Aortic Aneurysm}

\section*{Definition and Overview}
The aorta is the largest artery in the body, originating at the heart and passing through the chest and abdomen before dividing to supply the legs. An abdominal aortic aneurysm (AAA) is a permanent, localised widening of the aorta within the abdominal cavity, conventionally defined as a diameter of at least 3 cm. Most AAAs develop silently, grow slowly over years, and produce no symptoms. The principal danger is rupture: tearing of the weakened wall leads to catastrophic internal bleeding and a high risk of death. AAAs are detected either through screening programmes, incidentally on imaging performed for other reasons, or when they become symptomatic. Modern management offers surveillance for small aneurysms and elective repair for larger ones, substantially reducing the risk of rupture in those diagnosed in time.

\section*{Epidemiology}
AAAs are a disease of older adults, with prevalence rising steeply with age. Approximately 4--8\% of men over the age of 65 are found to have an AAA on ultrasound screening, while prevalence in women is several times lower. Risk is strongly associated with cumulative cigarette smoking, male sex, increasing age, a family history of aortic aneurysm, and the presence of hypertension and atherosclerotic vascular disease. In populations where smoking rates have declined and screening has been introduced, both the incidence of AAA and the rate of rupture have fallen. Rupture remains an important cause of sudden death, and in a substantial proportion of cases the aneurysm is discovered only at the point of rupture or at autopsy, underscoring the value of early detection.

\section*{Aetiology and Pathophysiology}
An AAA results from a complex process of structural failure of the arterial wall. The aorta's strength depends on its elastic and collagen fibres within the media; in aneurysm formation these are progressively degraded by chronic inflammation and an imbalance between matrix-degrading enzymes and their inhibitors. The wall consequently thins, weakens, and dilates under the constant pressure of arterial blood, and the law of Laplace ensures that as the diameter increases, the wall tension rises further, accelerating growth and eventual rupture.

\begin{itemize}
    \item \textbf{Smoking:} the strongest modifiable risk factor, promoting both the development and the expansion of AAAs
    \item \textbf{Age:} arterial wall integrity naturally declines with advancing years
    \item \textbf{Hypertension:} elevated blood pressure increases wall stress and the rate of expansion
    \item \textbf{Atherosclerosis and dyslipidaemia:} associated cardiovascular risk factors commonly coexist
    \item \textbf{Family history:} a first-degree relative with an AAA approximately doubles an individual's risk
    \item \textbf{Male sex:} men are affected four to five times more often than women
    \item \textbf{Connective tissue disorders:} Marfan syndrome, Ehlers--Danlos syndrome, and related conditions predispose to aneurysms, often at younger ages
    \item \textbf{Inflammation and infection:} inflammatory aortitis and rarely mycotic (infected) aneurysms
\end{itemize}

\section*{Clinical Features}
The majority of AAAs are asymptomatic and discovered incidentally or by screening. When symptoms do occur, they usually indicate that the aneurysm is large, expanding rapidly, or rupturing, and always warrant urgent assessment.

\begin{itemize}
    \item A pulsatile mass may be palpable in the abdomen on examination, although this is unreliable, especially in larger patients
    \item Vague abdominal, back, or flank discomfort in some individuals with large aneurysms
    \item Pressure effects on adjacent structures are uncommon but can cause symptoms
    \item A rapidly enlarging or leaking aneurysm produces new or worsening central abdominal or back pain
    \item Rupture classically presents with sudden, severe abdominal or back pain, collapse, pallor, sweating, and features of hypovolaemic shock
    \item Occasionally a ruptured AAA causes groin or thigh pain if bleeding tracks along the retroperitoneum
\end{itemize}

\section*{Differential Diagnosis}
The symptoms of a symptomatic or ruptured AAA overlap with many other causes of abdominal and back pain.

\begin{itemize}
    \item Musculoskeletal back pain and acute muscle strain
    \item Renal colic from kidney stones
    \item Acute pancreatitis, gallstone disease, and biliary colic
    \item Peptic ulcer disease and other causes of upper gastrointestinal pain
    \item Diverticulitis and other colonic conditions
    \item Aortic dissection, in which the layers of the aortic wall separate
    \item Perforated viscus or other causes of the acute abdomen
    \item Other causes of collapse and shock, including myocardial infarction
\end{itemize}

\section*{When to Seek Medical Care}
People with risk factors, including men aged 65 and over and those with a family history of AAA, should discuss screening with their doctor. Anyone who notices a new, persistent pulsating sensation in the abdomen should seek review. Sudden abdominal or back pain in a person at risk of an AAA must be treated as an emergency.

\textbf{Call 000 (or local emergency services) for:}
\begin{itemize}
    \item Sudden, severe, or tearing pain in the abdomen, back, or flank
    \item Collapse, faintness, or dizziness accompanied by abdominal or back pain
    \item Abdominal pain combined with pallor, sweating, and weakness
    \item Sudden severe chest pain or breathlessness, which may indicate aortic dissection or rupture
\end{itemize}

\section*{Diagnosis and Investigations}
Diagnosis is essentially radiological, because clinical examination cannot reliably detect or exclude an AAA.

\begin{itemize}
    \item \textbf{Abdominal ultrasound:} the first-line investigation and the mainstay of screening; it measures aortic diameter accurately and is safe, quick, and inexpensive
    \item \textbf{Computed tomography (CT) angiography:} provides detailed anatomy of the aneurysm and branch vessels and is essential for planning elective or emergency repair
    \item \textbf{Magnetic resonance imaging (MRI):} useful in selected cases where CT is contraindicated
    \item \textbf{Physical examination:} may reveal a pulsatile epigastric mass, but a normal examination does not exclude an AAA
    \item \textbf{Screening:} men aged 65 years are offered a one-off abdominal ultrasound in Australia; earlier testing is appropriate for smokers or those with a first-degree relative with an AAA
\end{itemize}

\section*{Management}
Management depends on aneurysm size, rate of growth, symptom status, and the patient's overall fitness.

\begin{itemize}
    \item \textbf{Surveillance:} small aneurysms below the repair threshold are monitored with interval ultrasound, typically every 6--12 months, with the interval guided by diameter
    \item \textbf{Cardiovascular risk factor control:} smoking cessation, blood pressure and cholesterol management, and a heart-healthy lifestyle slow aneurysm expansion
    \item \textbf{Elective repair:} indicated when the aneurysm is large (typically 5.5 cm in men, 5.0 cm in women), expands rapidly, or becomes symptomatic
    \item \textbf{Endovascular aneurysm repair (EVAR):} a minimally invasive technique in which a stent-graft is delivered through the femoral artery and deployed inside the aneurysm
    \item \textbf{Open surgical repair:} the aneurysm is excluded and replaced with a synthetic graft through an abdominal incision
    \item \textbf{Emergency repair:} for rupture, with either approach depending on anatomy and availability, although survival after rupture remains poor
\end{itemize}

\section*{Complications}
\begin{itemize}
    \item Rupture with massive retroperitoneal or intraperitoneal haemorrhage, the major and frequently fatal complication
    \item Intraluminal thrombus, which may rarely embolise and block arteries to the legs, kidneys, or bowel
    \item Pressure on adjacent structures, causing back pain or, uncommonly, ureteric or bowel symptoms
    \item Perioperative complications of repair, including bleeding, graft infection, kidney injury, bowel ischaemia, and myocardial infarction
    \item Endoleak after EVAR, in which blood continues to fill the aneurysm sac and may require further intervention
    \item Erectile dysfunction can follow open surgery as a result of pelvic nerve injury
\end{itemize}

\section*{Prognosis and Outcomes}
The natural history of an untreated AAA is of gradual expansion, with the annual risk of rupture rising sharply as diameter exceeds 5.5 cm. With surveillance, most small aneurysms never reach the point of requiring intervention. Elective repair is highly effective: operative mortality is low in fit patients, and after successful repair the risk of future rupture is effectively eliminated. Outcomes after EVAR are excellent in the short term, although some patients require secondary procedures; open repair offers durable long-term results. Ruptured AAA remains a devastating event, with overall mortality of 50--80\% even with emergency surgery. For this reason, screening, risk factor control, and timely elective repair are the cornerstones of improving outcomes.

\section*{Prevention and Self-Care}
\begin{itemize}
    \item Do not smoke, and seek cessation support if needed; smoking is the single most important modifiable risk factor
    \item Keep blood pressure and cholesterol within target ranges and attend routine checks
    \item Eat a balanced, heart-healthy diet and remain physically active
    \item Maintain a healthy weight and limit alcohol intake
    \item Ask about ultrasound screening if you are a man aged 65 or older
    \item Inform your doctor of any family history of aortic aneurysm or dissection so that earlier screening can be arranged
\end{itemize}

\section*{Sources and Further Reading}
The following organisations provide authoritative, patient-focused information on abdominal aortic aneurysm, screening, and treatment.

\begin{itemize}
    \item NHS. \textit{Abdominal aortic aneurysm: symptoms, screening, and treatment.} https://www.nhs.uk/conditions/abdominal-aortic-aneurysm/
    \item Mayo Clinic. \textit{Abdominal aortic aneurysm overview and management.} https://www.mayoclinic.org/diseases-conditions/abdominal-aortic-aneurysm/
    \item MSD Manual. \textit{Abdominal aortic aneurysms: professional reference.} https://www.msdmanuals.com/
    \item MedlinePlus. \textit{Abdominal aortic aneurysm health information.} https://medlineplus.gov/aorticaneurysm.html
    \item Centers for Disease Control and Prevention. \textit{Abdominal aortic aneurysm: risk factors and screening.} https://www.cdc.gov/
    \item NICE. \textit{Abdominal aortic aneurysm diagnosis and management.} https://www.nice.org.uk/guidance/ng156
\end{itemize}
